\scanpage{26}
\uline{Proposition}\quad \unsure{Soient|Soit} $\Sigma$ une $j$-Stratégie, $x_0\in C$ \deleted{position} gagnante pour $\Sigma$, $X=(x_0,x_1,\ldots,x_n\ldots)$ une partie (d'origine $x_0$) jouée suivant $\Sigma$, alors toutes les positions $x_i$ sont gagnantes pour $\Sigma$.

Soit en effet $Y=(y_0=x_i,y_1,\ldots,y_m,\ldots)$ une partie \inserted{maximale} jouée suivant $\Sigma$, \deleted{\unsure{et}} alors $XY$ est jouée suivant $\Sigma$ et est d'origine $x_0$ gagnante pour $\Sigma$, donc \deleted{est gagnée par $j$ donc $Y$ est finie} donc $Y$ est gagnée par $j$.

\uline{Th} On a \deleted{\unsure{rappel}}
\[
\begin{aligned}
a)\quad&G(j)\cap C(j)=\{x_0\in C(j)\mid R(x_0)\cap G(j)\ne\varnothing\}\\
b)\quad&G(j)\cap\bar C(j)=\{x_0\in\bar C(j)\mid R(x_0)\subset G(j)\}
\end{aligned}
\]
[et pour mémoire
\[
G(j)\cap C_0=\unsure{C_0|G_0}(j)
\]
].

\uline{Dém}\quad a) On suppose $x_0\in C(j)$.

\deleted{a)} 1) $\subset$ Soit $x_0\in G(j)\deleted{\cap C(j)}$, \inserted{soit} $\Sigma$ une stratégie pour $j$ gagnante pour $x_0$, considérons $x_1\in\Sigma(x_0)$, alors $(x_0,x_1)$ est une $\Sigma$-partie, donc par la prop. précédente, $x_1$ est gagnant pour $\Sigma$, donc $j$-gagnant, donc $x_1\in R(x)\cap G(j)$, donc $R(x)\cap G(j)\ne\varnothing$.

2) $\supset$ Soit $x_1\in R(x)\cap G(j)$. Prouvons que $x_0\in G(j)$. \deleted{\unsure{On défi|On déf}} Soit $\Sigma_1$ une $j$-stratégie $x_1$-gagnante. Si $x_0$ est gagnant pour $\Sigma_1$, on a $x_0\in G(j)$, OK. Sinon, définissons une $j$-stratégie $\Sigma_0$ $x_0$-gagnante ainsi :
\[
\begin{aligned}
\Sigma_0(x_0)&=\{x_1\}\\
\Sigma_0(x)&=\Sigma_1(x)\quad\text{si }x\ne x_0
\end{aligned}
\]
% archive foliation: 25
